\begingroup\sloppy
\begin{references}
\item National Cancer Institute Surveillance, Epidemiology, and End Results Program. Cancer stat facts: common cancer sites. Published 2020. https://seer.cancer.gov/statfacts/html/common.html.
\item Tilman D, Clark M. Global diets link environmental sustainability and human health. Nature. 2014;515(7528):518-522.
\item Bacchetti P. Current sample size conventions: flaws, harms, and alternatives. BMC Med. 2010;8:17.
\item Newman T, Kohn M. Evidence-Based Diagnosis: An Introduction to Clinical Epidemiology. 2nd ed. Cambridge University Press; 2020:292-295.
\item Hernan MA, Sauer BC, Hernandez-Diaz S, Platt R, Shrier I. Specifying a target trial prevents immortal time bias and other self-inflicted injuries in observational analyses. J Clin Epidemiol. 2016;79:70-75.
\item Newman TB. Antibiotic treatment for inpatient asthma exacerbations: what have we learned? JAMA Intern Med. 2021;181(4):427-428.
\item Newman TB. Possible immortal time bias in study of antibiotic treatment and outcomes in patients hospitalized for asthma. JAMA Intern Med. 2021;181(4):568-569.
\item Prasad V, Jena AB. Prespecified falsification end points: can they validate true observational associations? JAMA. 2013;309(3):241-242.
\item McEvoy SP, Stevenson MR, McCartt AT, et al. Role of mobile phones in motor vehicle crashes resulting in hospital attendance: a case-crossover study. BMJ. 2005;331(7514):428.
\item Magruder JT, Elahi D, Andersen DK. Diabetes and pancreatic cancer: chicken or egg? Pancreas. 2011;40(3):339-351.
\item Huxley R, Ansary-Moghaddam A, de Gonzalez Berrington A, Barzi F, Woodward M. Type-II diabetes and pancreatic cancer: a meta-analysis of 36 studies. Br J Cancer. 2005;92(11):2076-2083.
\sourcepage{311}
\item Bosetti C, Rosato V, Polesel J, et al. Diabetes mellitus and cancer risk in a network of case-control studies. Nutr Cancer. 2012;64(5):643-651.
\item Love HJ, Sulikowski D. Of meat and men: sex differences in implicit and explicit attitudes toward meat. Front Psychol. 2018;9:559.
\item Abotchie PN, Vernon SW, Du XL. Gender differences in colorectal cancer incidence in the United States, 1975-2006. J Womens Health (Larchmt). 2012;21(4):393-400.
\item Song M, Chan AT, Sun J. Influence of the gut microbiome, diet, and environment on risk of colorectal cancer. Gastroenterology. 2020;158(2):322-340.
\item Maconochie N, Doyle P, Carson C. Infertility among male UK veterans of the 1990-1 Gulf war: reproductive cohort study. BMJ. 2004;329(7459):196-201.
\item Trabert B, Sigurdson AJ, Sweeney AM, Strom SS, McGlynn KA. Marijuana use and testicular germ cell tumors. Cancer. 2011;117(4):848-853.
\item Bor J. Capitalizing on natural experiments to improve our understanding of population health. Am J Public Health. 2016;106(8):1388-1389.
\item Lofgren RP, Gottlieb D, Williams RA, Rich EC. Post-call transfer of resident responsibility: its effect on patient care [see comments]. J Gen Intern Med. 1990;5(6):501-505.
\item Bell CM, Redelmeier DA. Mortality among patients admitted to hospitals on weekends as compared with weekdays. N Engl J Med. 2001;345(9):663-668.
\item Davey Smith G, Ebrahim S. ‘Mendelian randomization’: can genetic epidemiology contribute to understanding environmental determinants of disease? Int J Epidemiol. 2003;32(1):1-22.
\item Ware JJ, van den Bree MB, Munafo MR. Association of the CHRNA5-A3-B4 gene cluster with heaviness of smoking: a meta-analysis. Nicotine Tob Res. 2011;13(12):1167-1175.
\item Chen LS, Hung RJ, Baker T, et al. CHRNA5 risk variant predicts delayed smoking cessation and earlier lung cancer diagnosis-a meta-analysis. J Natl Cancer Inst. 2015;107(5):djv100.
\item Zhou W, Zhu W, Tong X, et al. CHRNA5 rs16969968 polymorphism is associated with lung cancer risk: a meta-analysis. Clin Respir J. 2020;14(6):505-513.
\item Pillai SG, Ge D, Zhu G, et al. A genome-wide association study in chronic obstructive pulmonary disease (COPD): identification of two major susceptibility loci. PLoS Genet. 2009;5(3):e1000421.
\item Yang Q, Millard LAC, Davey Smith G. Proxy gene-by-environment Mendelian randomization study confirms a causal effect of maternal smoking on offspring birthweight, but little evidence of long-term influences on offspring health. Int J Epidemiol. 2020;49(4):1207-1218.
\item Doaei S, Hajiesmaeil M, Aminifard A, Mosavi-Jarrahi SA, Akbari ME, Gholamalizadeh M. Effects of gene polymorphisms of metabolic enzymes on the association between red and processed meat consumption and the development of colon cancer; a literature review. J Nutr Sci. 2018;7:e26.
\item Wang H, Iwasaki M, Haiman CA, et al. Interaction between red meat intake and NAT2 genotype in increasing the risk of colorectal cancer in Japanese and African Americans. PLoS One. 2015;10(12):e0144955.
\item Greenland S. An introduction to instrumental variables for epidemiologists. Int J Epidemiol. 2000;29(6):1102.
\item Rassen JA, Schneeweiss S, Glynn RJ, Mittleman MA, Brookhart MA. Instrumental variable analysis for estimation of treatment effects with dichotomous outcomes. Am J Epidemiol. 2009;169(3):273-284.
\item Bekelis K, Gottlieb DJ, Su Y, et al. Comparison of clipping and coiling in elderly patients with unruptured cerebral aneurysms. J Neurosurg. 2017;126(3):811-818.
\item Bor J, Moscoe E, Mutevedzi P, Newell ML, Barnighausen T. Regression discontinuity designs in epidemiology: causal inference without randomized trials. Epidemiology. 2014;25(5):729-737.
\item Hilton Boon M, Craig P, Thomson H, Campbell M, Moore L. Regression discontinuity designs in health: a systematic review. Epidemiology. 2021;32(1):87-93.
\item Moscoe E, Bor J, Barnighausen T. Regression discontinuity designs are underutilized in medicine, epidemiology, and public health: a review of current and best practice. J Clin Epidemiol. 2015;68(2):122-133.
\item Leyland AH, Ouedraogo S, Nam J, et al. Evaluation of Health in Pregnancy grants in Scotland: A Natural Experiment Using Routine Data. Public Health Research; 2017.
\item VanderWeele T. On the distinction between interaction and effect modification. Epidemiology. 2009;20(6): 863-871.
\item Braitman LE, Rosenbaum PR. Rare outcomes, common treatments: analytic strategies using propensity scores. Ann Intern Med. 2002;137(8):693-695.
\item Gum PA, Thamilarasan M, Watanabe J, Blackstone EH, Lauer MS. Aspirin use and all-cause mortality among patients being evaluated for known or suspected coronary artery disease: a propensity analysis. JAMA. 2001;286(10):1187-1194.
\item Bruns SB, Ioannidis JP. p-Curve and p-hacking in observational research. PLoS One. 2016;11(2):e0149144.
\item Newman TB, Bernzweig JA, Takayama JI, Finch SA, Wasserman RC, Pantell RH. Urine testing and urinary tract infections in febrile infants seen in office settings: the Pediatric Research in Office Settings’ Febrile Infant Study. Arch Pediatr Adolesc Med. 2002;156(1):44-54.
\item Katz MH. Multivariable Analysis: A Practical Guide for Clinicians. 2nd ed. Cambridge University Press; 2006:xv, 192 p.
\end{references}
\endgroup
